% !TEX root = ../notes.tex

\documentclass[../notes.tex]{subfiles}

\begin{document}

\section{September 10}
Today, we discuss how to do algebra with $\infty$-categories. Today's notes were copied (with gratitude) from Yutong Chen.

\subsection{Presentable Categories}
We are going to deal with size issues today, so let's recall a little set theory.
\begin{definition}[regular]
	Fix a cardinal $\kappa$. Then a set $S$ is \textit{$\kappa$-small} if and only if $\left|S\right|<\kappa$; in other words, there is an injection $S\into\kappa$. A cardinal $\kappa$ is \textit{regular} if and only if every $\kappa$-small union of $\kappa$-small sets is still $\kappa$-small.
\end{definition}
\begin{example}
	All successor cardinals of $\aleph_0$ are regular. For example, $\aleph_0$ is regular because a finite union of finite sets is finite. Similarly, $\aleph_1$ is regular because a countable union of countable sets is countable.
\end{example}
\begin{remark}
	Let $S$ be a set whose cardinality is a regular cardinal. Then the well-ordering principle allows us to order the elements of $S$ as $\{s_\alpha\}_{\alpha\in\kappa}$, where $\kappa$ has a natural ordering. This ordering produces a filtration
	\[S=\bigcup_{\alpha\in\kappa}S_\alpha,\]
	where each $S_\alpha$ is $\kappa$-small, and any $\kappa$-small subset $T$ lives in $S_\alpha$ for some $\alpha$. One can show this latter claim by a quick instance of transfinite induction on the size of $T$.
\end{remark}
We now switch from sets to categories.
\begin{definition}[filtered]
	Fix a regular cardinal $\kappa$. A category $I$ is \textit{$\kappa$-filtered} if and only if every functor $D\to I$ from a $\kappa$-small category $D$ factors through the category $D^\Delta$, which is the category $D$ together with a terminal object.
\end{definition}
\begin{remark}
	Roughly speaking, we are asking for the functor $F\colon D\to I$ to admit an object $y\in I$ such that each $x\in D$ has a morphism $Fx\to y$. (This is only a rough approximation because the maps $Fx\to y$ should also commute with each other.) For example, $I$ is $\kappa$-filtered if $I$ is closed under $\kappa$-small colimits.
\end{remark}
\begin{definition}[compact]
	Fix a regular cardinal $\kappa$. Fix an $\infty$-category $\mc C$. An object $X\in\mc C$ is \textit{$\kappa$-compact} if and only if $\op{Hom}_{\mc C}(X,-)$ commutes with all $\kappa$-filtered colimits.
\end{definition}
\begin{remark}
	Intuitively, $\kappa$-compact means that the object is presented by fewer than $\kappa$ elements.
\end{remark}
\begin{example}
	Let $\mc C$ be the (homotopy) category of abelian groups. Then a finitely presented abelian group is $\aleph_0$-compact.
\end{example}
\begin{proposition}
	Fix a $\kappa$ cardinal $\kappa$. If an $\infty$-cateogry is $\kappa$-compactly generated, then $\kappa$-filtered colimits commute with $\kappa$-small products.
\end{proposition}
\begin{proof}
	Omitted. For example, we have not defined what it means to be compactly generated.
\end{proof}
Here is the main definition of this subsection.
\begin{definition}[presentable]
	Fix a regular cardinal $\kappa$. A category $\mc C$ is \textit{$\kappa$-presentable} if and only if it has all $\kappa$-small colimits, and there is a set $S\subseteq\mc C$ of $\kappa$-compact objects for which each $X\in\mc C$ is a colimit of elements in $S$.
\end{definition}
\begin{remark}
	Any object in a $\kappa$-presentable category $\mc C$ can be written as a $\kappa$-filtered colimit of compact objects. Indeed, write $X=\colim_{i\in I}K_i$ where $K_i\in S$ for each $i$. Then the key is that $I=\colim_{\alpha\in\kappa} I_\alpha$ as in our set-theoretic discussion, so we are done upon writing
	\[X=\colim_{i\in I}K_i=\colim_{\alpha\in\kappa}\colim_{i\in I_\alpha}K_i.\]
	One now checks that each object $\colim_{i\in I_\alpha}K_i$ is $\kappa$-compact.
\end{remark}

\subsection{A Category of Categories}
Here are some basic structural properties.
\begin{notation}
	Let $\mc C$ be a category, and let $\kappa$ be a regular cardinal. Then we let $\mc C^\kappa$ denote the full subcategory of $\kappa$-compact objects.
\end{notation}
\begin{theorem}
	Let $\mc C$ be a $\kappa$-presentable category. Let $\op{Ind}_\kappa\mc C^\kappa$ be the category obtained by formally adding $\kappa$-filtered colimits. Then the natural functor
	\[\op{Ind}_\kappa\mc C^\kappa\to\mc C\]
	given by taking the formal colimits to the actual colimits is an equivalence.
\end{theorem}
\begin{proof}
	We omit the proof. It is formal.	
\end{proof}
\begin{proposition} \label{prop:c-kappa-is-essentially-small}
	Let $\mc C$ be a $\kappa$-presentable category. Then $\mc C^\kappa$ is essentially $\kappa$-small.
\end{proposition}
\begin{proof}
	For any $\kappa$-compact object $K\in\mc C$, we may write
	\[K=\colim_{i\in I}S_i\]
	for some $S_i\in S$, where $I$ is $\kappa$-filtered. By compactness, the natural map $K\to\colim_iS_i$ factors through $S_i$ for some $i$, so $K$ is a retract of $S_i$. Thus, we may enumerate the isomorphism classes in $\mc C^\kappa$ by enumerating the retracts of the objects in $S$, which is a $\kappa$-small collection.
\end{proof}
With a smallness condition, we can now build a category of categories.
\begin{notation}
	Fix a regular cardinal $\kappa$. Let $\op{Pr}^L_\kappa$ denote the category of $\kappa$-presentable categories, where morphisms are functors which preserve colimits and $\kappa$-compact objects. Furthermore, we define $\op{Pr}^L$ to be the ``union'' of all $\op{Pr}^L_\kappa$, which is the category of categories which are $\kappa$-presentable for some $\kappa$, where morphisms are functors which preserve colimits.
\end{notation}
\begin{remark}
	The union $\op{Pr}^L$ has some size issues with its $\op{Hom}$-sets: for example,
	\[\op{Hom}_{\op{Pr}^L}(\mathrm{Ani},\mc C)=\mc C\]
	for any $\mc C\in\op{Pr}^L$. (Indeed, a functor out of $\mathrm{Ani}$ which preserves colimits is uniquely determined by where it sends the point.)
\end{remark}
\begin{notation}
	Fix a regular cardinal $\kappa$. Let $\op{Rex}_\kappa$ be the category of all $\kappa$-small categories which admit $\kappa$-small colimits and retracts, where morphisms are functors preserving $\kappa$-small colimits.
\end{notation}
\begin{theorem}[Gabriel--Ulmer] \label{thm:gab-ulm}
	The natural functors $(-)^\kappa\colon{\op{Pr}_\kappa^L}\to{\op{Rex}_\kappa}$ and $\op{Ind}_\kappa\colon{\op{Rex}_\kappa}\to\op{Pr}_\kappa^L$ are equivalences.
\end{theorem}
\begin{proof}[Sketch]
	We only discuss one of the inverse checks: we will show that $\left(\op{Ind}_\kappa\mc C\right)^\kappa=\mc C$. (The other inverse check was stated but not proved previously.) Well, start with some $\mc C\in\op{Rex}_\kappa$. On one hand, one can check that $\mc C\subseteq\op{Ind}_\kappa\mc C$ maps to the $\kappa$-compact objects, due to the formal nature of the colimits we added. On the other hand, for any object $X\in\op{Ind}_\kappa\mc C$ which is $\kappa$-compact, we may write $X=\colim_{\alpha\in I}D_\alpha$ for some $\kappa$-filtered $I$ by construction of $\op{Ind}_\kappa\mc C$; then $X$ is the retract of one of the objects $D_\alpha$, so $X\in\mc C$ by definition of $\op{Rex}_\kappa$.
\end{proof}
\begin{remark}
	\Cref{thm:gab-ulm} tells us that the $\op{Hom}$-sets in $\op{Pr}^L_\kappa$ do not have the size issues that $\op{Pr}^L$ does: indeed,
	\[\op{Hom}_{\op{Pr}^L_\kappa}(\mc C,\mc D)=\op{Hom}_{\op{Rex}_\kappa}(\mc C^\kappa,\mc D^\kappa),\]
	and the target is $\kappa$-small. For example, each of $\mc C^\kappa$ and $\mc D^\kappa$ are essentially $\kappa$-small by \Cref{prop:c-kappa-is-essentially-small}.
\end{remark}
We will view $\op{Pr}^L_\kappa$ as our ``category of categories.'' For example, one desirable feature of this definition is that it is (essentially) an element of itself.
\begin{theorem}
	The category $\op{Pr}^L_\kappa$ is $\kappa$-presentable. We may write this mathematical statement as $\mathrm{Pr}^L_\kappa\in\mathrm{Pr}^L_\kappa$.
\end{theorem}
\begin{proof}[Sketch]
	By \Cref{thm:gab-ulm}, it is enough to show that $\op{Rex}_\kappa$ is $\kappa$-presentable. We will study the forgetful functor $U\colon\op{Rex}_\kappa\to\mathrm{Cat}$. For example, it reflects isomorphisms. Also, \Cref{thm:gab-ulm} implies that $F\colon\mathrm{Cat}_\infty\to\op{Rex}_\kappa$ defined by
	\[F(S)\coloneqq\op{PSh}(S;\mathrm{Ani})^\kappa\]
	is a left adjoint for $F$. Here, $F$ is better thought of as the corresponding functor $\op{PSh}\colon\mathrm{Cat}_\infty\to\op{Pr}^L_\kappa$. I don't want to think about the adjunction checks.

	Now, $\mathrm{Cat}_\infty$ is generated by the $\kappa$-compact object $\Delta^n$. Also, the functor $F$ sends $\kappa$-compact objects to $\kappa$-compact objects. Thus, $\{F\Delta^n\}_n$ are $\kappa$-compact objects, and we would like them to generate. To this end, define $T\coloneqq FG$, so there is a natural map $T\to1$ given by the adjunction. It turns out that the natural map
	\[T^2\mc C\to T\mc C\]
	has co-equalizer $\mc C$, which completes the proof by writing the relevant objects in $\mathrm{Cat}_\infty$ as colimits of $\Delta^n$s.
\end{proof}

\subsection{Algebra with Higher Categories}
Another reason why presentability is a desirable hypothesis is the presence of an adjoint functor theorem.
\begin{theorem}[Lurie]
	Let $F\colon\mc C\to\mc D$ be a morphism in $\op{Pr}^L_\kappa$.
	\begin{listalph}
		\item Then $F$ admits a right adjoint if and only if it commutes with small colimits.
		\item Then $F$ admits a left adjoint if and only if $F$ commutes with small limits and is accessible, meaning that $F$ commutes with $\kappa$-filtered colimits for some regular cardinal $\kappa$.
	\end{listalph}
\end{theorem}
\begin{remark}
	Roughly speaking, there is an asymmetry in (a) and (b) because presentability is about generation by colimits. If one instead defined the corresponding category $\op{Pr}^R_\kappa$ (whose morphisms preserve accessible functors), then we would expect the asymmetry to happen in the other direction.
\end{remark}
\begin{remark}
	It turns out that $\mathrm{Pr}^R_\kappa\cong\op{Pr}^R$, and the forgetful functor $\mathrm{Pr}^R\to\mathrm{Cat}_\infty$ preserves small limits. One can use this to show that $\op{Pr}^L$ admits small colimits.
\end{remark}
\begin{example}
	Coproducts in $\op{Pr}^L$ are (na\"ive) products. Products are also na\"ive products. 
\end{example}
\begin{example}
	Products in $\op{Pr}^L_\kappa$ are $\op{Ind}$-saturations of na\"ive products.
\end{example}
Products and coproducts being the same means that we will have to work harder to build a tensor product.
\begin{definition}[Lurie tensor product]
	Fix $\mc C,\mc D\in\op{Pr}^L_\kappa$. Then we define $\mc C\otimes\mc D$ by universal property: we have
	\[\op{Hom}_{\op{Pr}^L_\kappa}(\mc C\otimes\mc D,-)\cong\op{Bilin}_{\op{Pr}^L_\kappa}(\mc C\times\mc D,-).\]
	Here, ``$\op{Bilin}$'' refers to the functors linear in both coordinates.
\end{definition}
\begin{remark}
	One can construct $\mc C\otimes\mc D$ via the adjoint functor theorem. Namely, the functor $\mc C\otimes-$ is an adjoint to the corresponding bilinear functor.
\end{remark}
% \begin{example}
% 	Suppose that $\mc C$ is presented as $\langle S\rangle\left[R^{-1}\right]$, and $\mc D$ is presented as $\langle T\rangle\left[Q^{-1}\right]$. Then $\mc C\otimes D$ can be checked to be presented as $\langle S\times T\rangle\left[(R\times Q)^{-1}\right]$.
% \end{example}

\end{document}